\documentclass[pdf, unicode, 12pt, a4paper, oneside, fleqn]{article}

\makeatletter
\usepackage[utf8]{inputenc}
\usepackage[T2B]{fontenc}
\usepackage[english,russian]{babel}

\usepackage{geometry}
\geometry{
  a4paper,
  left=2.54cm,
  right=2.54cm,
  top=2.54cm,
  bottom=2.54cm,
  headheight=0mm,
  headsep=0mm,
  footnotesep=0mm,
  footskip=0mm,
}

\usepackage{nopageno}

\usepackage{tempora}
\renewcommand{\ttdefault}{cmtt}

\usepackage{titlesec}
\titleformat*{\section}{\large\bfseries}
\titlespacing\section{0pt}{12pt plus 4pt minus 2pt}{0pt plus 2pt minus 2pt}
\setcounter{secnumdepth}{0}

\usepackage{enumitem}

\frenchspacing

\usepackage{amsmath}
\usepackage{amssymb}
\usepackage{hyperref}
\usepackage{longtable}
\usepackage[table]{xcolor}
\usepackage{array}
\usepackage{color}
\usepackage{xcolor}

\usepackage{hyperref}

\newcommand{\MYhref}[3][blue]{\href{#2}{\color{#1}{#3}}}%

\usepackage{listings}
\usepackage{alltt}
\usepackage{csquotes}
\DeclareQuoteStyle{russian}
	{\guillemotleft}{\guillemotright}[0.025em]
	{\quotedblbase}{\textquotedblleft}
\ExecuteQuoteOptions{style=russian}

\usepackage{graphicx}

\lstset{tabsize=2,
	breaklines,
	columns=fullflexible,
	flexiblecolumns,
	numbers=left,
	numberstyle={\footnotesize},
	extendedchars,
	inputencoding=utf8}

\def\@xobeysp{ }
\def\verbatim@processline{\hspace{1.2cm}\raggedright\the\verbatim@line\par}

\parindent=0pt
\parskip=3pt

\definecolor{lightgray}{gray}{0.9}

\renewcommand{\thesubsection}{\arabic{subsection}}

\lstdefinestyle{customc}{
  belowcaptionskip=1\baselineskip,
  breaklines=true,
  frame=L,
  xleftmargin=\parindent,
  language=C,
  showstringspaces=false,
  basicstyle=\footnotesize\ttfamily,
  keywordstyle=\bfseries\color{green!40!black},
  commentstyle=\itshape\color{gray},
  identifierstyle=\color{black},
  stringstyle=\color{blue},
}

\lstdefinestyle{customasm}{
  belowcaptionskip=1\baselineskip,
  frame=L,
  xleftmargin=\parindent,
  language=[x86masm]Assembler,
  basicstyle=\footnotesize\ttfamily,
  commentstyle=\itshape\color{purple!40!black},
}

\lstset{escapechar=@,style=customc}
\makeatother

\usepackage{pgfplots}
\pgfplotsset{compat=1.16}


\lstset{keepspaces=true}
\lstset{literate=
  {а}{{\cyra}}1 {б}{{\cyrb}}1 {в}{{\cyrv}}1 {г}{{\cyrg}}1 {д}{{\cyrd}}1 {е}{{\cyre}}1
  {ё}{{\cyryo}}1 {ж}{{\cyrzh}}1 {з}{{\cyrz}}1 {и}{{\cyri}}1 {й}{{\cyrishrt}}1 {к}{{\cyrk}}1
  {л}{{\cyrl}}1 {м}{{\cyrm}}1 {н}{{\cyrn}}1 {о}{{\cyro}}1 {п}{{\cyrp}}1 {р}{{\cyrr}}1
  {с}{{\cyrs}}1 {т}{{\cyrt}}1 {у}{{\cyru}}1 {ф}{{\cyrf}}1 {х}{{\cyrh}}1 {ц}{{\cyrc}}1
  {ч}{{\cyrch}}1 {ш}{{\cyrsh}}1 {щ}{{\cyrshch}}1 {ъ}{{\cyrhrdsn}}1 {ы}{{\cyrery}}1 {ь}{{\cyrsftsn}}1
  {э}{{\cyrerev}}1 {ю}{{\cyryu}}1 {я}{{\cyrya}}1 {А}{{\CYRA}}1 {Б}{{\CYRB}}1 {В}{{\CYRV}}1
  {Г}{{\CYRG}}1 {Д}{{\CYRD}}1 {Е}{{\CYRE}}1 {Ё}{{\CYRYO}}1 {Ж}{{\CYRZH}}1 {З}{{\CYRZ}}1
  {И}{{\CYRI}}1 {Й}{{\CYRISHRT}}1 {К}{{\CYRK}}1 {Л}{{\CYRL}}1 {М}{{\CYRM}}1 {Н}{{\CYRN}}1
  {О}{{\CYRO}}1 {П}{{\CYRP}}1 {Р}{{\CYRR}}1 {С}{{\CYRS}}1 {Т}{{\CYRT}}1 {У}{{\CYRU}}1
  {Ф}{{\CYRF}}1 {Х}{{\CYRH}}1 {Ц}{{\CYRC}}1 {Ч}{{\CYRCH}}1 {Ш}{{\CYRSH}}1 {Щ}{{\CYRSHCH}}1
  {Ъ}{{\CYRHRDSN}}1 {Ы}{{\CYRERY}}1 {Ь}{{\CYRSFTSN}}1 {Э}{{\CYREREV}}1 {Ю}{{\CYRYU}}1 {Я}{{\CYRYA}}1
  {—}{{---}}1 {«}{{<<}}1 {»}{{>>}}1 {№}{{\textnumero}}1
}


\begin{document}

\begin{titlepage}
\begin{center}
\large
МОСКОВСКИЙ АВИАЦИОННЫЙ ИНСТИТУТ\\ (НАЦИОНАЛЬНЫЙ ИССЛЕДОВАТЕЛЬСКИЙ УНИВЕРСИТЕТ)

\vspace{20pt}

Институт №8 <<Компьютерные науки и прикладная математика>>

Кафедра 806 <<Вычислительная математика и программирование>>
\end{center}

\vspace{60pt}

\begin{center}
\bfseries
\large
Лабораторная работа \textnumero 3 \\по курсу \enquote{Операционные системы}

\vspace{54pt}

\end{center}

\vfill

\begin{flushright}
\large
\begin{tabular}{rl}
Выполнил: & Яковлев А. С. \\
Группа: & М8О-203БВ-25 \\
Преподаватель: & Е.\,С.~Миронов \\
\end{tabular}
\end{flushright}

\vspace{92pt}

\begin{center}
\large
Москва, \the\year
\end{center}
\end{titlepage}

\pagebreak

\section{Условие}
{\bfseries Цель работы:} приобретение практических навыков в
\begin{itemize}[nosep]
    \item управлении потоками в ОС;
    \item обеспечении синхронизации между потоками.
\end{itemize}

{\bfseries Задание:} составить программу на языке Си, обрабатывающую данные
в многопоточном режиме. При обработке использовать стандартные средства создания
потоков операционной системы (Windows/Unix). Ограничение максимального количества
потоков, работающих в один момент времени, должно быть задано ключом запуска программы.
Продемонстрировать количество потоков, используемое программой, с помощью стандартных
средств операционной системы. Привести исследование зависимости ускорения и эффективности
алгоритма от входных данных и количества потоков и объяснить полученные результаты.

{\bfseries Вариант:} 9. Рассчитать детерминант матрицы (используя определение детерминанта).

\section{Метод решения}

\subsection*{Определитель по определению}

Определитель квадратной матрицы $A = (a_{ij})$ порядка $n$ по определению равен
\[
    \det A = \sum_{\sigma \in S_n} \operatorname{sgn}(\sigma)\, a_{0,\sigma(0)}\, a_{1,\sigma(1)} \cdots a_{n-1,\sigma(n-1)},
\]
где сумма берётся по всем $n!$ перестановкам $\sigma$ множества $\{0, \dots, n-1\}$.
Каждое слагаемое~--- произведение $n$ элементов, взятых по одному из каждой строки
и каждого столбца. Знак перестановки определяется числом инверсий~--- пар $i < j$,
для которых $\sigma(i) > \sigma(j)$: при чётном числе инверсий $\operatorname{sgn}(\sigma) = +1$,
при нечётном~--- $-1$.

Перестановки перебираются рекурсивно: функция \texttt{go}
выбирает для очередной строки свободный столбец, помечает его занятым, домножает
накопленное произведение на выбранный элемент и переходит к следующей строке.
Число инверсий накапливается по ходу: при выборе столбца $c$ в строке $k$ новых
инверсий столько, сколько уже занятых столбцов больше $c$ (они стоят в перестановке
раньше и правее). Когда заполнены все строки, произведение возвращается со знаком,
определяемым чётностью числа инверсий.

\subsection*{Распараллеливание}

Слагаемые суммы независимы, поэтому работу можно разделить между потоками без обмена
данными по ходу вычислений. Вся сумма разбита на $n^2$ задач: задача с номером $t$
фиксирует столбцы первых двух строк $c_0 = \lfloor t / n \rfloor$, $c_1 = t \bmod n$
и перебирает все перестановки с таким началом. Задачи с $c_0 = c_1$ пусты, остальные
$n(n-1)$ задач содержат по $(n-2)!$ слагаемых. Каждая перестановка попадает ровно в одну
задачу, поэтому сумма результатов всех задач равна определителю. Разбиение по двум строкам,
а не по одной, даёт мелкие задачи (при $n = 12$~--- 132 непустые), которые равномернее
распределяются по потокам.

Задачи раздаются динамически: общий счётчик \texttt{next\_task} указывает номер следующей
свободной задачи. Освободившийся поток под мьютексом берёт номер и увеличивает счётчик.
Так потоки, которые планировщик ОС выполняет медленнее (например, на энергоэффективных
ядрах), берут меньше задач и не задерживают остальных.

Каждый поток накапливает сумму своих задач в локальной переменной и по окончании работы
один раз прибавляет её к общему результату под тем же мьютексом.

\subsection*{Ограничение числа потоков}

Ключ \texttt{-t P} задаёт общее число одновременно работающих потоков, включая главный.
Программа создаёт $P - 1$ потоков функцией \texttt{pthread\_create}, а главный поток
выполняет ту же функцию \texttt{worker} наравне с ними. При $P = 1$ дополнительных
потоков не создаётся. Затем главный поток ожидает завершения остальных (\texttt{pthread\_join}).

\section{Описание программы}

\subsection*{Структура проекта}

\begin{center}
\begin{tabular}{|l|p{11cm}|}
\hline
\textbf{Файл} & \textbf{Содержимое} \\ \hline
\texttt{CMakeLists.txt} & сборка: библиотека \texttt{detlib} и программа \texttt{det}, подключение pthreads \\ \hline
\texttt{src/main.c} & разбор ключей запуска, получение матрицы, замер времени, вывод результата \\ \hline
\texttt{src/matrix.c}, \texttt{.h} & чтение матрицы из файла, генерация случайной матрицы, печать \\ \hline
\texttt{src/det.c}, \texttt{.h} & вычисление определителя, управление потоками и синхронизация \\ \hline
\texttt{tests/m1..m4.txt} & тестовые матрицы с известным определителем \\ \hline
\end{tabular}
\end{center}

Матрица $n \times n$ хранится в одномерном массиве \texttt{double} по строкам:
элемент $a_{ij}$ находится по индексу \texttt{i * n + j}.

\subsection*{Основные функции}

\begin{itemize}
    \item \texttt{double *matrix\_read(const char *filename, int *n)}~--- читает матрицу из файла, возвращает массив или \texttt{NULL} при ошибке.
    \item \texttt{double *matrix\_random(int n)}~--- создаёт матрицу из случайных целых
          чисел от $-9$ до $9$. Генератор инициализируется константой (\texttt{srand(1)}),
          поэтому при одинаковом $n$ матрица всегда одна и та же, что делает замеры
          сопоставимыми.
    \item \texttt{double go(int row, int *used, double prod, int inv)}~--- рекурсивно
          перебирает перестановки для строк начиная с \texttt{row}, возвращает сумму
          соответствующих слагаемых.
    \item \texttt{double do\_task(int t)}~--- фиксирует столбцы первых двух строк
          по номеру задачи и вызывает \texttt{go} со строки~2.
    \item \texttt{void *worker(void *arg)}~--- функция потока: в цикле берёт задачи
          из общей очереди, накапливает локальную сумму, в конце добавляет её к результату.
    \item \texttt{double det\_calc(double *a, int n, int threads)}~--- инициализирует общие
          данные, создаёт $P-1$ потоков, выполняет \texttt{worker} в главном потоке,
          дожидается остальных и возвращает определитель.
\end{itemize}

\subsection*{Ключи запуска}

\begin{tabular}{ll}
\texttt{-t P} & число потоков (обязательный, от 1 до 256) \\
\texttt{-f F} & прочитать матрицу из файла \texttt{F} \\
\texttt{-n N} & сгенерировать случайную матрицу $N \times N$ ($N \le 15$) \\
\texttt{-p}   & вывести матрицу \\
\end{tabular}

Программа печатает свой PID, определитель и время вычисления.

\subsection*{Использованные источники}
{\sloppy
\begin{enumerate}[nosep]
    \item Таненбаум~Э., Бос~Х. Современные операционные системы.~--- СПб.:
\end{enumerate}
}

\section{Результаты}

\subsection*{Конфигурация}

\begin{tabular}{ll}
Процессор: & M3 Pro \\
Ядер всего: & 11 \\
ОС: & MacOS \\
\end{tabular}

\subsection*{Проверка корректности}

Программа запускалась на матрицах с заранее известным определителем при разном числе
потоков. Ожидаемые значения.

\begin{center}
\begin{tabular}{|l|r|r|r|r|r|}
\hline
\textbf{Матрица} & \textbf{Ожидается} & $P=1$ & $P=2$ & $P=4$ & $P=8$ \\ \hline
\texttt{m1.txt} ($1\times1$) & 7      & 7 & 7 & 7 & 7 \\ \hline
\texttt{m2.txt} ($2\times2$) & $-14$  & -14 & -14 & -14 & -14 \\ \hline
\texttt{m3.txt} ($3\times3$) & $-306$ & -306 & -306 & -306 & -306 \\ \hline
\texttt{m4.txt} ($4\times4$) & 30     & 30 & 30 & 30 & 30 \\ \hline
\end{tabular}
\end{center}

Результат не зависит от числа потоков, что подтверждает отсутствие гонок данных:
каждая задача вычисляется ровно один раз, а общий результат изменяется только под мьютексом.

\subsection*{Демонстрация числа потоков}

Программа была запущена командой \texttt{./det -t 4 -n 13}, во время вычислений
число её потоков просмотрено стандартными средствами ОС. Получено 4 потока: главный
и три созданных.

\begin{lstlisting}[language={},numbers=none,escapechar={},columns=fixed,basicstyle=\footnotesize\ttfamily]
artem@MacBook-Pro-Artem-3 lab3 % ps -M -p 45502
USER    PID   TT   %CPU STAT PRI     STIME     UTIME COMMAND
artem 45502 s000  100.0 R    31T   0:00.01   0:04.50 ./cmake-build-debug/det -t 4 -n 13
      45502       100.0 R    31T   0:00.01   0:04.49
      45502       100.0 R    31T   0:00.01   0:04.49
      45502       100.0 R    31T   0:00.01   0:04.49
\end{lstlisting}

\subsection*{Исследование ускорения и эффективности}

Замеры выполнялись на случайных целочисленных матрицах порядка $n = 8 \dots 12$
(одинаковых для всех $P$) при числе потоков $P = 1, 2, 4, 8, 16$. Каждый замер
повторялся 5 раз, в таблицу занесена медиана.

Ускорение $S(P) = T_1 / T_P$, эффективность $E(P) = S(P) / P \cdot 100\%$,
где $T_1$~--- время при одном потоке, $T_P$~--- при $P$ потоках.

\begin{center}
Время $T_P$, с

\begin{tabular}{|c|r|r|r|r|r|}
\hline
$P$ & $n=8$ & $n=9$ & $n=10$ & $n=11$ & $n=12$ \\ \hline
1  & 0.000714 & 0.006422 & 0.0665 & 0.7614 & 9.696 \\ \hline
2  & 0.000399 & 0.003333 & 0.0341 & 0.3911 & 4.849 \\ \hline
4  & 0.000236 & 0.001750 & 0.0179 & 0.2024 & 2.487 \\ \hline
8  & 0.000204 & 0.001220 & 0.0118 & 0.1396 & 1.791 \\ \hline
16 & 0.000216 & 0.001140 & 0.0107 & 0.1193 & 1.558 \\ \hline
\end{tabular}

\vspace{8pt}
Ускорение $S(P)$

\begin{tabular}{|c|r|r|r|r|r|}
\hline
$P$ & $n=8$ & $n=9$ & $n=10$ & $n=11$ & $n=12$ \\ \hline
1  & 1.00 & 1.00 & 1.00 & 1.00 & 1.00 \\ \hline
2  & 1.79 & 1.93 & 1.95 & 1.95 & 2.00 \\ \hline
4  & 3.03 & 3.67 & 3.72 & 3.76 & 3.90 \\ \hline
8  & 3.50 & 5.26 & 5.66 & 5.45 & 5.41 \\ \hline
16 & 3.31 & 5.63 & 6.22 & 6.38 & 6.22 \\ \hline
\end{tabular}

\vspace{8pt}
Эффективность $E(P)$, \%

\begin{tabular}{|c|r|r|r|r|r|}
\hline
$P$ & $n=8$ & $n=9$ & $n=10$ & $n=11$ & $n=12$ \\ \hline
1  & 100.0 & 100.0 & 100.0 & 100.0 & 100.0 \\ \hline
2  & 89.5 & 96.3 & 97.6 & 97.3 & 100.0 \\ \hline
4  & 75.6 & 91.7 & 93.0 & 94.0 & 97.5 \\ \hline
8  & 43.8 & 65.8 & 70.8 & 68.2 & 67.7 \\ \hline
16 & 20.7 & 35.2 & 38.9 & 39.9 & 38.9 \\ \hline
\end{tabular}
\end{center}

\begin{center}
\begin{tikzpicture}
\begin{axis}[
  width=0.68\textwidth, height=7cm,
  xmode=log, log basis x=2,
  xtick={1,2,4,8,16}, xticklabels={1,2,4,8,16},
  xmin=0.9, xmax=18, ymin=0, ymax=8,
  xlabel={Число потоков $P$}, ylabel={Ускорение $S$},
  grid=major, legend pos=outer north east, legend cell align=left,
  unbounded coords=jump, cycle list name=color list,
  every axis plot/.append style={thick, mark=*},
]
  \addplot[dashed, gray, domain=1:16, samples=2, mark=none] {x};
  \addlegendentry{идеал $S=P$}
  %                     P=1      P=2      P=4      P=8      P=16
    \addplot coordinates {(1,1.00) (2,1.79) (4,3.03) (8,3.50) (16,3.31)}; \addlegendentry{$n=8$}
  \addplot coordinates {(1,1.00) (2,1.93) (4,3.67) (8,5.26) (16,5.63)}; \addlegendentry{$n=9$}
  \addplot coordinates {(1,1.00) (2,1.95) (4,3.72) (8,5.66) (16,6.22)}; \addlegendentry{$n=10$}
  \addplot coordinates {(1,1.00) (2,1.95) (4,3.76) (8,5.45) (16,6.38)}; \addlegendentry{$n=11$}
  \addplot coordinates {(1,1.00) (2,2.00) (4,3.90) (8,5.41) (16,6.22)}; \addlegendentry{$n=12$}
\end{axis}
\end{tikzpicture}

Рис. 1. Зависимость ускорения от числа потоков

\vspace{12pt}

\begin{tikzpicture}
\begin{axis}[
  width=0.68\textwidth, height=7cm,
  xmode=log, log basis x=2,
  xtick={1,2,4,8,16}, xticklabels={1,2,4,8,16},
  xmin=0.9, xmax=18, ymin=0, ymax=110,
  xlabel={Число потоков $P$}, ylabel={Эффективность $E$, \%},
  grid=major, legend pos=outer north east, legend cell align=left,
  unbounded coords=jump, cycle list name=color list,
  every axis plot/.append style={thick, mark=*},
]
  \addplot[dashed, gray, domain=1:16, samples=2, mark=none] {100};
  \addlegendentry{идеал $E=100\%$}
  %                       P=1        P=2      P=4      P=8      P=16
  \addplot coordinates {(1,100.0) (2,89.5) (4,75.6) (8,43.8) (16,20.7)}; \addlegendentry{$n=8$}
  \addplot coordinates {(1,100.0) (2,96.3) (4,91.7) (8,65.8) (16,35.2)}; \addlegendentry{$n=9$}
  \addplot coordinates {(1,100.0) (2,97.6) (4,93.0) (8,70.8) (16,38.9)}; \addlegendentry{$n=10$}
  \addplot coordinates {(1,100.0) (2,97.3) (4,94.0) (8,68.2) (16,39.9)}; \addlegendentry{$n=11$}
  \addplot coordinates {(1,100.0) (2,100.0) (4,97.5) (8,67.7) (16,38.9)}; \addlegendentry{$n=12$}
\end{axis}
\end{tikzpicture}

Рис. 2. Зависимость эффективности от числа потоков
\end{center}

\section{Выводы}

В ходе работы была написана программа на языке Си, вычисляющая определитель матрицы
по определению~--- как сумму $n!$ слагаемых по всем перестановкам~--- в многопоточном
режиме с использованием POSIX Threads. Работа разбивается на $n^2$ независимых задач,
которые потоки берут из общей очереди; доступ к общим данным (счётчику задач и результату)
защищён мьютексом. Число потоков ограничивается ключом запуска с учётом главного потока,
что подтверждено средствами ОС. Корректность проверена на матрицах с известным
определителем: результат не зависит от числа потоков.

Исследование на процессоре Apple M3 Pro (5 производительных и 6 энергоэффективных ядер)
показало следующее.

\begin{enumerate}
    \item \textbf{До 4 потоков ускорение почти линейное.} При $n = 12$ получено
          $S(2) = 2{,}00$ и $S(4) = 3{,}90$ (эффективность 100\% и 97,5\%). Задачи независимы,
          мьютекс захватывается лишь на мгновение при выдаче задачи, а матрица (около 1~КБ)
          целиком помещается в кэш каждого ядра, поэтому задача ограничена только
          вычислениями. Все потоки при этом работают на производительных ядрах.

    \item \textbf{Неоднородные ядра.} При $P = 8$ ускорение растёт заметно слабее
          ($S = 5{,}41$, $E = 67{,}7\%$ при $n = 12$): производительных ядер только 5,
          и остальные потоки выполняются на более медленных энергоэффективных ядрах.

    \item \textbf{Потоков больше, чем ядер.} При $P = 16$, что больше 11 ядер, ускорение
          выходит на плато: $S \approx 6{,}2$--$6{,}4$ для $n \ge 10$. Это близко к
          предельной производительности всех ядер процессора. Лишние потоки лишь делят
          процессорное время, поэтому эффективность падает до 35--40\%.

    \item \textbf{Малые матрицы.} При $n = 8$ (40\,320 слагаемых) всё вычисление занимает
          около 0,7~мс, и затраты на создание и завершение потоков становятся заметными.
          Поэтому эффективность ниже уже при $P = 4$ (75,6\%), ускорение не превышает 3,5,
          а при $P = 16$ программа даже немного замедляется по сравнению с $P = 8$
          (0,216~мс против 0,204~мс). По закону Амдала доля последовательной части
          (создание потоков, ожидание их завершения) здесь велика. С ростом $n$ эта доля
          уменьшается, и эффективность растёт.

\end{enumerate}

Таким образом, задача хорошо распараллеливается: независимые слагаемые и минимальная
синхронизация дают почти линейное ускорение, пока потоки работают на производительных
ядрах, а его предел определяется числом и типом ядер процессора и накладными расходами
на потоки при малом объёме работы.

\pagebreak

\end{document}